\documentclass[10pt,a4paper]{amsart}
\usepackage[margin=19mm]{geometry}
\usepackage{amsmath,amssymb,mathtools}
\usepackage[scr=rsfs,cal=euler]{mathalfa}
\usepackage{xfrac}
\usepackage[unicode,hidelinks,bookmarks=false]{hyperref}
\newcommand{\ie}{i.e.\ }
\newcommand{\cf}{cf.\ }
\newcommand{\resp}{respectively\ }
\newcommand{\Subsection}{\subsection}
\providecommand{\nobreakdash}{\nobreak\hbox{-}}
\setlength{\emergencystretch}{2em}
\multlinegap0pt
\usepackage{bm}
\usepackage{mathtools}
\usepackage{amssymb}
\newtheorem{lemma}{Lemma}
\title{A Note on the Point-Clothoid Distance Algorithm}
\author{Haibin Ye, Hao Ge, Gong Cheng}
\date{Source: arXiv:2609.10179v1 (CC BY 4.0)}
\begin{document}
\maketitle

\noindent\emph{Source and licence.} A Note on the Point-Clothoid Distance Algorithm. Version arXiv:2609.10179v1 (CC BY 4.0), distributed by arXiv under the Creative Commons Attribution 4.0 International licence, http://creativecommons.org/licenses/by/4.0/, as recorded in the arXiv licence metadata for this exact version and shown on the abstract page https://arxiv.org/abs/2609.10179v1. Full licence notice: \texttt{licenses/arxiv-2609.10179-CC-BY-4.0.txt}.\par\smallskip

\noindent\textit{Editorial scope.} Counted unit: the article's lemma lem:convexity (source0.tex lines 291--294). The article's complete original proof of it is delivered with it (source0.tex lines 296--319, one \texttt{proof} environment of 24 lines). No mathematics is written, repaired or re-derived: the statement and the proof are the article's own lines, in the article's own notation, and no retained line is substituted or re-flowed.  No further source-local context is needed: every cross-reference of every retained line resolves inside the two delivered spans.  Nothing was omitted from any retained span, so the package carries no omission record.  No retained line contains a citation, so the wrapper carries no bibliography and no bibliography entry.  No retained line contains a figure environment, an image-inclusion command or drawing code, so no image is delivered and the package declares no asset dependency.  Editorial formatting only, in addition: an AMS \texttt{amsart} preamble that restates the article's own declarations and macros used by the retained lines, namely the environments \texttt{lemma} on separate counters, together with the standard packages those lines need.  The retained environments print in this document as Lemma 1 in order of appearance, whereas the article prints the same items with the numbering of its own full document.  The complete native source of the article as distributed in the arXiv e-print for this exact version is bundled unchanged as \texttt{sources/209882/source0.tex}.
\begin{lemma}
  \label{lem:convexity}
  A function $f$ constructed as above is strictly convex on $(0,M)$.
\end{lemma}

\begin{proof}
  For every segment--frame pair, the same calculation follows after shifting
  or reversing the local tangent angle. It therefore suffices to display it
  for $\Gamma_2$ in the frame
  $[\boldsymbol{e}(s_{\pi/2});-\mathbf{N}(s_{\pi/2}),
  \mathbf{T}(s_{\pi/2})]$. Let
  $(x_\gamma(s),y_\gamma(s))$, $s\in[s_{\pi/2},s_\pi]$, denote its
  coordinates. Since $\theta(s)=ks^2/2+k_0s$ and
  $\kappa(s)=ks+k_0$,
  \begin{equation}
    \label{eq:derivative of para_gamma_2}
    x_\gamma'(s)=\frac{k\sin\theta(s)}{\kappa(s)^2},
    \qquad
    y_\gamma'(s)=-\frac{k\cos\theta(s)}{\kappa(s)^2}.
  \end{equation}
  For $s\in(s_{\pi/2},s_\pi)$, $x_\gamma'(s)>0$, and hence
  \begin{equation}
    \label{derivative of f}
    f'(x)=\frac{y_\gamma'(s)}{x_\gamma'(s)}=-\cot\theta(s),
    \qquad
    f''(x)=\frac{\kappa(s)^3}{k\sin^3\theta(s)}>0.
  \end{equation}
  Thus $f$ is strictly convex on $(0,M)$.
\end{proof}

\end{document}
